\documentclass{article}
\usepackage[T2A]{fontenc}
\usepackage[english, russian]{babel}
\usepackage{mathtools, amssymb, amsfonts, amsthm}
\usepackage{tikz-cd}
\usetikzlibrary{babel}
\usepackage{enumitem}

\newtheorem*{theorem}{Теорема}
\newtheorem*{corollary}{Следствие}
\newtheorem*{lemma}{Лемма}
\newtheorem*{statement}{Утверждение}
\newtheorem*{proposition}{Предложение}
\theoremstyle{definition}
\newtheorem*{definition}{Определение}
\theoremstyle{remark}
\newtheorem*{properties}{Свойства}
\newtheorem*{reminder}{Напоминание}
\newtheorem*{remark}{Замечание}
\newtheorem*{example}{Пример}
\newtheorem{problem}{Задача}

\DeclareMathOperator{\Hom}{Hom}
\DeclareMathOperator{\coker}{coker}
\DeclareMathOperator{\Ob}{Ob}
\DeclareMathOperator{\Arr}{Arr}
\newcommand{\Cat}{\mathcal{C}}
\newcommand{\A}{\mathcal{A}}

\newenvironment{enum}[2][\parindent]{\begin{enumerate}[nosep, wide, labelindent = #1, label=#2]}{\end{enumerate}}
\newenvironment{cd}{\begin{equation*}\begin{tikzcd}}{\end{tikzcd}\end{equation*}}
\renewcommand{\thesection}{Глава \arabic{section}}
\renewcommand{\thesubsection}{Глава \arabic{section}.\arabic{subsection}}

\title{}
\author{}
\date{}

\begin{document}
\maketitle

\section{Триангулированные категории}

\setcounter{subsection}{-1}
\subsection{Абелевы категории}

\begin{definition}
    Аддитивной категорией называется категория $\Cat$ со структурой абелевой группы на каждом $\Hom(X, Y)$ такой, что:
    \begin{enum}{\arabic*)}
        \item $f \circ (g + h) = f \circ g + f \circ h$.
        \item Существует $0$ --- инициальный и терминальный объект.
        \item Существует копроизведение = произведение двух объектов.
    \end{enum}
\end{definition}

\begin{definition}
    $\mathcal{C}$ --- адитивная категория, $f \in \Hom(X, Y)$, $g \in \Hom(K, X)$ называется ядром, если $f \circ g = 0$ и для любого $\tilde g$: $f \circ \tilde{g} = 0$ существует единственный $h$, что диаграмма коммутирует.
\end{definition}

\begin{cd}
    K \rar["g"] & X \rar["f"] & Y \\
    \tilde{K} \ar[u, "h"] \ar[ur, "\tilde{g}"] 
\end{cd}

\begin{definition}
    Аддитивная категория предабелева, если любой морфизм имеет ядро и коядро.
\end{definition}

$\Cat$ --- предабелева, $f \in \Hom(X, Y)$.
\begin{cd}
\ker \rar & X \dar \rar["f"] \ar[dr, "\tilde{f}"] & Y \dar \rar & \coker \\
& \coker(\ker) \rar["\tilde{\tilde{f}}"] & \ker(\coker)
\end{cd}

\begin{definition}
    Если в предабелевой категории $\Cat$ $\tilde{\tilde{f}}$ --- изоморфизм, $\Cat$ абелева.
\end{definition}

\begin{theorem}[Фрейда--Митчелла]
    Любая абелева категория вкладывается как полная подкатегория в $A-mod$ для некоторого $A$ аддитивным точным функтором.
\end{theorem}

$\A$ --- абелева катугория, $C(\A)$ --- категория комплексов, она тоже абелева.

\begin{theorem}[О длинной гомомлогической последовательности]
    Если
    \begin{cd}
        0 \rar & A \rar & B \rar & C \rar & 0
    \end{cd}
    --- коротка точная последовательность комплексов, существует длинная последовательность
    \begin{cd}
        \ldots \rar & H^n(A) \rar & H^n(B) \rar & H^n(C) \rar["\delta_n"] & H^{n + 1}(A) \rar & \ldots
    \end{cd}
    $H^n$ --- набор функторов. 
\end{theorem}

В новой категории

\begin{enum}{\arabic*)}
    \item Для короткой точной последовательности
    \begin{cd}
        0 \rar & A \rar & B \rar & C \rar & 0
    \end{cd}
    возникает $h: C \rightarrow A[1]$ и последовательность
    \begin{cd}
        0 \rar & A \rar["f"] & B \rar["g"] & C \rar["h"] & A[1] \rar["{-f[1]}"] & \ldots
    \end{cd}
    \item Любой $s \in \Hom_{\Cat(A)}(X, Y)$ в новой категории $s = f$, где $f$ --- моно(В частности возникает длинная последовательность). Итого, любой $f$ включается в "выделенный треугольник".
\end{enum}

\subsection{Категории с надстройкой}

\begin{definition}
    $\Cat$ --- аддитивная категория, $S$ --- аддитивный эндофунктор.
    \begin{cd}
        X \rar["f"] & Y \rar["g"] & Z \rar["h"] & S(X)
    \end{cd}
    --- $S$ последовательность. Они образвуют аддитивную категорию
    \begin{cd}
        X \dar["f"] \rar & Y \dar["g"] \rar & Z \dar["h"] \rar & S(X) \dar["S(f)"] \\
        X' \rar & Y' \rar & Z' \rar & S(X')
    \end{cd}
\end{definition}

Пусть задан класс $\Delta$ $S$-последовательностей.

$(\Cat, S, \Delta)$ называется категорией с надстройкой, если выполнены аксиомы:
\begin{enum}{Tr \arabic*}
    \setcounter{enumi}{-1}
    \item $\Delta$ замкнут относительно изоморфизмов.
    \item $\Delta$ содержит все $S$-последовательности вида $0 \rightarrow A \xrightarrow{id} A \rightarrow 0$.
    \item $Y \xrightarrow{g} Z \xrightarrow{h} S(X) \xrightarrow{-Sf} S(Y) \in \Delta$.
    \item Если есть две последовательности и коммутативный квадрат
    \begin{cd}
        X \dar["u"] \rar["f"] & Y \dar["v"] \\
        X' \rar["f'"] & Y'
    \end{cd}
    то он достраивается до морфизма $S$-последовательностей.
    \item Любой квадрат 
    \begin{cd}
        A \dar[equal] \rar & B \dar \\
        A \rar & C
    \end{cd}
    достраивается до диаграммы
    \begin{cd}
    A \dar[equal] \rar & B \dar \rar & D \dar \rar & S(A) \dar[equal]\\
    A \rar & C \dar \rar & E \dar \rar & S(A) \dar["S(f)"]\\
    & K \dar \rar[equal] & K \dar \rar & S(B) \\
    & S(B) \rar & S(D)
    \end{cd}
\end{enum}

\begin{lemma}
    Если
    \begin{cd}
        X \rar["u"] & Y \rar["v"] & Z \rar["w"] & SX \in \Delta,
    \end{cd}
    то $vu = 0$.
\end{lemma}

\begin{proof}
    \begin{cd}
    X \rar["id"] \dar["id"] & X \dar["u"] \rar & 0 \dar[dashed] \rar & SX \dar[dashed] \\
    X \rar["u"] & Y \rar["v"] & Z \rar["w"] & SX
    \end{cd}
\end{proof}

\begin{lemma}
    \begin{cd}
        X \rar["u"] & Y \rar["v"] & Z \rar["w"] & SX
    \end{cd}
    $\in \Delta$, тогда для любого $T$.
    \begin{cd}
    \Hom(Z, T) \rar & \Hom(Y, T) \rar & \Hom(X, T)
    \end{cd}
    точна
\end{lemma}
\begin{proof}
    \begin{cd}
        X \rar["u"] \dar & Y \rar["v"] \dar["f"] & Z \rar["w"] \dar[dashed] & SX \dar[dashed]\\
        0 \rar & T \rar["id"] & T \rar & 0
    \end{cd}
\end{proof}

\begin{lemma}
    \begin{cd}
        X \rar \dar["f"] & Y \rar \dar["g"] & Z \rar \dar["h"] & SX \dar["Sf"] \in \Delta \\
        X' \rar & Y' \rar & Z' \rar & SX' \in \Delta
    \end{cd}
    Если $f, g$ --- изоморфизмы, то и $h$ --- изоморфизм.
\end{lemma}
\begin{proof}
    Сдвинем треугольники применим лемму выше
    \begin{cd}
        \Hom(SY, T) \rar & \Hom(SX, T) \rar & \Hom(Z, T) \rar & \Hom(Y, T) \rar & \Hom(X, T) \\
        \Hom(SY', T) \uar \rar & \Hom(SX', T) \rar \uar & \Hom(Z', T) \uar["h^*"] \rar & \Hom(Y', T) \uar \rar & \Hom(X', T) \uar
    \end{cd}
    Т. е. для любого $T$ по 5-лемме $h^*$ --- изоморфизм, так что по лемме Йонеды $h$ --- изоморфизм.
\end{proof}

\begin{definition}
    $f \in \Hom(X, Y)$, $(Z = C(f), g, h)$ --- конус морфизма, если
    \begin{cd}
        X \rar["f"] & Y \rar["g"] & Z \rar["h"] & SX \in \Delta
    \end{cd}
    Из Tr 4 конус существует.
\end{definition}

\begin{corollary}
    $C(f)$ единственный с точностью до изоморфизма.
\end{corollary}

\begin{definition}
    \begin{cd}
        X \rar & Y \rar & Z \rar & SX
    \end{cd}
    --- предтреугольник, если последовательность
    \begin{cd}
        \Hom(SY, T) \rar & \Hom(SX, T) \rar & \Hom(Z, T) \rar & \Hom(Y, T) \rar & \Hom(X ,T)
    \end{cd}
    точна.
\end{definition}

\begin{lemma}
    Сумма $S$-последовательностей лежит в $\Delta$ тогда и только тогда, когда оба слагаемых лежат в $\Delta$.
\end{lemma}
\begin{proof}
    Пусть каждое из слагаемых лежит.
    \begin{cd}
    X \oplus X' \rar & Y \oplus Y' \rar & \tilde{Z} \rar & S(X \oplus X') \in \Delta
    \end{cd}
    Далее 5-лемма или диаграмма прямой суммы.
    
    В обратную сторону
    \begin{cd}
    X \dar \rar & Y \dar \rar & Z \dar[dashed] \rar & SX \dar[dashed] \\
    X \oplus X \dar' \rar & Y \oplus Y' \dar \rar & Z \oplus Z' \dar[dashed] \rar & S(X \oplus X') \dar[dashed] \in \Delta \\
    X \rar & Y \rar & \tilde{Z} \rar & SX \in \Delta
    \end{cd}
    Далее опять 5-лемма.
\end{proof}

\begin{lemma}
    \begin{cd}
        X \rar["0"] & Y \rar & Z \rar & SX \in Delta
    \end{cd}
    Тогда он изоморфен
    \begin{cd}
        X \rar["0"] & Y \rar & Y \oplus SX \rar & SX
    \end{cd}
\end{lemma}
\begin{proof}
    Вторая принадлежит $\Delta$ как сумма, так что она изоморфна.
\end{proof}

\begin{definition}
    Триангулированная категория --- категория с надстройкой, в которой $S$ эквивалентность.

    Предтриангулированная --- триангулированная категория без октаэдра, но $\Hom$-ы достраиваются.
\end{definition}

\begin{properties} \leavevmode
    \begin{enum}{\arabic*)}
        \item Если
        \begin{cd}
            B \rar & C \rar & S(A) \rar["-Sf"] & S(B) \in \Delta,
        \end{cd}
        то и
        \begin{cd}
            A \rar["f"] & B \rar & C \rar & S(A) \in \Delta
        \end{cd}
        \item Можно достроить любой квадрат(необязательно левый)
        \item 
        \begin{cd}
        \Hom(T, X) \rar & \Hom(T, Y) \rar & \Hom(T, Z)
        \end{cd}
        точна
    \end{enum}
\end{properties}

\begin{lemma}
    $f \in \Hom(X, Y)$, если $f$ --- мономорфизм, $Y = X \oplus Z$. Аналогично с эпи.
\end{lemma}
\begin{proof}
    Докажем для эпи. $f$ --- эпи.
    \begin{cd}
        X \rar["f"] & Y \rar & Z \rar & SX
    \end{cd}
    $gf = 0$, значит $g = 0$. Тогда 
    \begin{cd}
        Y \rar["0"] \dar & Z \rar \dar & SX \rar \dar & SY \dar \in \Delta \\
        Y \rar & Z \rar & Z \oplus SY \rar & SY
    \end{cd}
\end{proof}

Есть отображение $\Ob(\Arr \Cat) \rightarrow \Ob \Cat$. Любая пара морфизмов достраивается, но это не функториально.

Факт. Если $\Cat$ триангулированная категория и $C(f)$ продолжается до функтора, то $\Cat$ --- полная триангулированная подкатегория в абелевой полупростой категория(с тривиальной структурой)

\begin{definition}
    Пусть $\Cat_1, \Cat_2$ триангулировнные. Триангулированным функтором из $\Cat_1$ в $\Cat_2$ называется $(F, R)$, $F: \Cat_1 \rightarrow \Cat_2$ --- аддитивный функтор. $R: S_{\Cat_2}F \leftarrow FS_{\Cat_1}$ --- изоморфизм, т. ч., если
    \begin{cd}
        X \rar & Y \rar & Z \rar & SX \in \Delta_{C_1},
    \end{cd}
    \begin{cd}
        FX \rar & FY \rar & FZ \rar & SFX \in \Delta_{C_2}
    \end{cd}
\end{definition}

\begin{definition}
    $\Cat_1$ триангулированная подкатегории, если $\Cat_1$ полная подкатегория, и $i: \Cat_1 \rightarrow \Cat_2$ --- триангулированный функтор.
\end{definition}

\begin{definition}
    Аддитивная категория $\Cat$ называется Карубиевой(идемпотентно полной), если все идемпотенты расщипляются, т. е. если $e \in \Hom(A, A)$, $e^2 = e$, то $A \cong A_1 \oplus A_2$ и диаграмма
    \begin{cd}
    A \rar \dar[equal] & A_1 \rar & A \dar[equal] \\
    A \ar[rr, "e"] & & A
    \end{cd}
    коммутирует.
\end{definition}

Факт. Любая триангулированная вкладывается в карубиеву.

Факт 2. $\Cat$ --- карубиева, конус функториален, тогда $\Cat$ полупростая абелева с тривиальной структурой.

\begin{proposition}
    Если в триангулированной категории идемпотенты расщепляются, и конус функториален, тогда любой треугольник сумма тривиальных.
\end{proposition}
\begin{proof}
    \begin{cd}
        X \rar["f"] & Y \rar["g"] \dar["g"] & Z \dar["c(g)"] \rar & SX \\
        0 \rar & Z \rar  \dar["c(g)"]& Z \rar \dar["c(g)"]& 0 \\
        0 \rar & Z \rar & Z \rar & 0
    \end{cd}
    т. е. $c(g)^2 = c(gc(g)) = c(g)$. Значит $C \simeq C_1 \oplus C_2$.
    \begin{cd}
        C \rar["p"] & C_1 \rar["i"] & C
    \end{cd}
    \begin{proposition}
        $pg$ --- эпи
    \end{proposition}
    \begin{proof}
    $upg = vpg$, тогда $c(upg) = c(vpg)$, значит $upc(g) = vpc(g)$, значит $up = vp$, значит $u = v$. 
    \end{proof}
    Но эпи расщипляется и любой $g \in \Hom(B, C)$ изоморфен $\begin{tikzcd} B \rar & C_1 \oplus C_2 \end{tikzcd}$. 
\end{proof}

Более общая теорема:

\begin{theorem}
    Копредел диаграммы --- коконус с универсальным свойством. Слабый копредел --- коконус для которого есть стрелка
\end{theorem}

\begin{statement}
Слабый копредел функториален, идемпотенты рсащепляются, тогда это предел.
\end{statement}

\section{Локализация категорий}

$\Cat$ --- категория, $M$ --- класс морфизмов. Локализация $C \xrightarrow{i} C(M^{-1})$ такой, что для любого $F: C \rightarrow D$, т. ч. $F(s)$  --- изоморфизмы, существуе $G$, т. ч. $F = G \circ I$.

\begin{definition}
    $M$ --- правый локализующий класс($\Cat$ --- аддитивная), если
    \begin{enum}{\arabic*)}
        \item $id_X \in M$
        \item $f, g \in M$, тогда $f \circ g \in M$.
        \item Если 
        \begin{cd}
        X \dar["u"] \rar["f"] & Z \\
        Y
        \end{cd}
        и $u \in M$, тогда есть 
        \begin{cd}
        X \dar["u"] \rar["f"] & Z \dar["v"] \\
        Y \rar["g"] & T
        \end{cd}
        где $v \in M$.
        \item 
        \begin{cd}
            Z \rar["u"] & X \rar["f"] & Y
        \end{cd}
        $u \in M$, $fu = 0$, тогда найдётся $v \in M$(в какой-то объект), т. ч. $vf = 0$.
    \end{enum}
\end{definition}

Построим $C[M^{-1}]$. Объекты те же. Определим $\Hom(X, Y)$.

$\tilde{\Hom}(X, Y) = \{X \xrightarrow{f} Z \xleftarrow{u} Y| u \in M\}$.

Зададим отношение эквивалентности
\begin{equation*}
    \begin{tikzcd}
        & Z \\
        X \ar[ur, "f"] & & Y \ar[ul, "u"]
    \end{tikzcd}
    \sim
    \begin{tikzcd}
        & Z' \\ 
        X \ar[ur, "f'"] & & Y \ar[ul, "u'"]
    \end{tikzcd}
    ,\text{ если }
\end{equation*}
\begin{cd}
    & T \\
    Z \ar[ur, "g"] & & Z' \ar[ul, "v"] \\
    X & & Y
\end{cd}
где $vu' \in M$.

$\Hom(X, Y)$ --- фактор класс

Композиция

\begin{cd}
& & S \\
Z & & T \\
X & & Y & & W
\end{cd}
(замкнули по 3).
\end{document}